\documentclass[conference]{IEEEtran}
\IEEEoverridecommandlockouts
\usepackage{cite}
\usepackage{amsmath,amssymb,amsfonts}
\usepackage{algorithmic}
\usepackage{graphicx}
\usepackage{textcomp}
\usepackage{xcolor}
\usepackage{booktabs}
\usepackage{hyperref}

\begin{document}

\title{Satellite Conjunction Risk Analysis with Collision Detection\\
\large A Deep-Learning-Based System for Satellite Conjunction Risk Analysis and Collision Detection with Comparative Monte Carlo Validation}

\author{\IEEEauthorblockN{[Candidate Name / Roll Number]}
\IEEEauthorblockA{\textit{[Department of Computer Science \& Engineering]} \\
\textit{[Institution / University Name]}\\
Supervised by: [Project Supervisor / Guide Name, Designation] \\
Academic Year: [Academic Year 2025–2026] \\
Target Repository: \texttt{E:\textbackslash DL\_fixed} (feature/operational-monte-carlo-surrogate)}}

\maketitle

\begin{abstract}
As the tracked resident space object (RSO) catalog exceeds 36,000 orbital bodies and commercial mega-constellations rapidly populate Low Earth Orbit (LEO), satellite operators face an unmanageable surge in conjunction alerts. Operational workflows rely on high-precision numerical orbit propagation and Monte Carlo covariance sampling to estimate collision probabilities ($P_c$), imposing severe computational bottlenecks and triggering chronic operator alert fatigue. In this project, deep learning is developed as the primary methodology for rapid, automated satellite conjunction risk analysis, miss distance prediction, and collision detection. We train deep sequence architectures---specifically Retentive Networks (RetNet) featuring multi-scale retention and Causal Standard Transformers---alongside tabular models to analyze historical Conjunction Data Message (CDM) sequences enriched with orbital physics from the Space-Track Satellite Catalog (SATCAT). Vectorized 1,000,000-particle B-plane Monte Carlo simulation is retained as an essential comparative benchmark to evaluate risk estimation fidelity, prediction behaviour, computational efficiency, probability differences, and operational trade-offs.

Under a rigorous 5-fold cross-validation protocol on 3,351 CDMs across 836 unique encounters (210 high-risk, 626 low-risk), tabular gradient boosting achieves the highest discriminative performance ($\text{ROC-AUC} = 0.9086 \pm 0.0226$, $\text{PR-AUC} = 0.8121 \pm 0.0439$). Among deep sequence architectures, Causal Transformer achieves $\text{ROC-AUC} = 0.8307 \pm 0.0369$, while RetNet achieves $\text{ROC-AUC} = 0.7506 \pm 0.0484$. Sequence analysis reveals that 59.5\% of events contain only 1 to 3 updates (median = 3.0), explaining why static orbital shell features dominate discrimination. Crucially, standardized single-core CPU profiling demonstrates that RetNet recurrent inference ($0.57\text{ ms}$) and Transformer inference ($0.78\text{ ms}$) achieve verified $80.0\times$ and $58.3\times$ speedups over 1,000,000-sample Monte Carlo quadrature ($45.28\text{ ms}$), while HistGradientBoosting ($1.10\text{ ms}$) delivers a $41.3\times$ speedup. In a triage pipeline, gradient boosting safely screens out 65.4\% of non-critical conjunctions, routing only flagged critical encounters to intensive Monte Carlo integration. We conclude that deep learning provides an effective, sub-millisecond screening surrogate that enhances space safety when integrated with numerical verification.
\end{abstract}

\begin{IEEEkeywords}
Satellite Conjunction Assessment, Collision Detection, Deep Learning Surrogates, Retentive Networks (RetNet), Causal Transformers, Monte Carlo Simulation, Space Situational Awareness (SSA), Conjunction Data Messages (CDMs).
\end{IEEEkeywords}

\section{Introduction}
The contemporary orbital environment is undergoing an unprecedented transformation. With more than 36,000 trackable objects exceeding 10 cm in Earth orbit and mega-constellations rapidly deploying thousands of satellites into dense Low Earth Orbit (LEO) shells, satellite operators receive hundreds of Conjunction Data Messages (CDMs) daily. High-fidelity numerical orbit propagation and Monte Carlo covariance sampling are computationally expensive, scaling quadratically ($O(N^2)$) across all-vs-all satellite pairings. 

In this research, deep learning is developed as the primary methodology for rapid satellite conjunction risk analysis and collision detection. We train deep neural architectures to analyze sequential historical CDM telemetry and forecast collision risk at sub-millisecond latency. Concurrently, vectorized 1,000,000-particle B-plane Monte Carlo simulation is retained as an essential comparative benchmark to evaluate risk estimation fidelity, latency speedup, and numerical differences.

\section{Problem Statement \& Objectives}
Let an orbital encounter $i$ be monitored over sequential CDMs: $\mathcal{S}_i = \{\mathbf{x}_1, \dots, \mathbf{x}_k\}$, with $1 \le k \le T_i \le 14$ and feature dimension $D = 14$. The task is formulated as classifying whether the encounter will escalate to high collision risk requiring mitigation maneuvers: $Y_i = 1$ if $\max_{t} P_c(t) \ge 5 \times 10^{-4}$, else $0$. The primary objective is to evaluate deep sequence models (RetNet and Causal Transformer) alongside tabular models and compare them directly against 1,000,000-sample Monte Carlo simulations.

\section{Feature Space \& Data Processing}
The dataset comprises 3,351 CDMs across 836 unique encounters (210 high-risk, 626 low-risk) fused with 70,000+ Space-Track SATCAT records. At each step, 14 physical channels are extracted: $\log_{10}(\text{miss})$, $\Delta \text{miss}$, range ratio, normalized step position, $\Delta \text{apogee}$, $\Delta \text{perigee}$, $\Delta \text{inclination}$, $\Delta \text{period}$, primary/secondary RCS, object classifications, and exclusion radii. For tabular models, 71 summary statistics are extracted (`summarize\_71`). To prevent padding leakage, sequences are dynamically packed via \texttt{pack\_real\_first}.

\section{Evaluated Models}
\begin{itemize}
    \item \textbf{RetNet Surrogate:} 3-layer Retentive Network using multi-scale matrix retention with decay parameters $\gamma \in \{0.90, 0.95, 0.98, 0.99\}$, SwiGLU feed-forward networks, and $d_{\text{model}} = 64$ (104,801 parameters). Supports parallel training and recurrent $O(1)$ streaming inference ($S_t = \gamma S_{t-1} + K_t^T V_t$).
    \item \textbf{Standard Transformer:} 3-layer Transformer encoder with strict causal upper-triangular masking and masked pooling over real tokens (155,201 parameters).
    \item \textbf{Classical Baselines:} Majority Class, 3-state Markov Chain, Autoregressive/ARIMA, Logistic Regression, Random Forest (200 trees), and HistGradientBoosting.
    \item \textbf{Monte Carlo Simulator:} Vectorized 1,000,000-sample 2D B-plane Gaussian random particle generator with 95\% Wilson confidence intervals.
\end{itemize}

\section{Experimental Results}
All models were benchmarked under identical 5-fold cross-validation on 836 events.

\begin{table}[htbp]
\caption{Model Comparison under Unified 5-Fold Cross-Validation}
\begin{center}
\begin{tabular}{lcccc}
\toprule
\textbf{Model Architecture} & \textbf{AUC-ROC} & \textbf{PR-AUC} & \textbf{Brier} & \textbf{Latency} \\
\midrule
Monte Carlo ($10^6$ particles) & Ref. & Ref. & N/A & 45.28 ms \\
Majority Class & $0.500 \pm 0.000$ & $0.251 \pm 0.001$ & 0.188 & N/B \\
Markov Chain & $0.578 \pm 0.020$ & $0.285 \pm 0.009$ & 0.183 & N/B \\
ARIMA / Autoregressive & $0.784 \pm 0.048$ & $0.635 \pm 0.083$ & 0.184 & N/B \\
Logistic Regression & $0.842 \pm 0.030$ & $0.670 \pm 0.065$ & 0.160 & 0.094 ms \\
Random Forest (200 trees) & $0.905 \pm 0.036$ & $0.796 \pm 0.068$ & 0.112 & 11.19 ms \\
\textbf{HistGradientBoosting} & $\mathbf{0.909 \pm 0.023}$ & $\mathbf{0.812 \pm 0.044}$ & \textbf{0.103} & 1.10 ms \\
Standard Transformer & $0.831 \pm 0.037$ & $0.644 \pm 0.063$ & 0.159 & 0.78 ms \\
RetNet (Proposed) & $0.751 \pm 0.048$ & $0.535 \pm 0.072$ & 0.182 & \textbf{0.57 ms} \\
\bottomrule
\end{tabular}
\end{center}
\end{table}

\section{Comparative Analysis \& Discussion}
Relative to 2D B-plane Monte Carlo numerical quadrature ($45.28\text{ ms}$ for $10^6$ particles), RetNet recurrent inference ($0.57\text{ ms}$) achieves an $80.0\times$ speedup, Transformer ($0.78\text{ ms}$) achieves $58.3\times$, and HistGradientBoosting ($1.10\text{ ms}$) achieves $41.3\times$. Tabular gradient boosting achieves highest discrimination ($\text{AUC} = 0.9086$) due to sequence length sparsity (median 3 CDMs, 59.5\% $\le 3$ updates), where static orbital shell features dominate. In a triage screening configuration, gradient boosting safely screens out 65.4\% of non-critical conjunctions while maintaining high safety recall.

\section{Limitations \& Flight Safety Disclaimers}
Public CDMs omit full 2x2 covariance matrices, requiring isotropic calibration under which 71\% of high-risk cases require hard-body radius inflation (median 3.8x) due to real-world along-track covariance elongation. Furthermore, the Collision Avoidance Maneuver (CAM) module in \texttt{orbit\_guard.py} implements a linearized along-track Clohessy-Wiltshire secular drift heuristic ($\Delta v = \text{gap} / (3 t)$) for demonstration purposes only, and is not certified for operational flight use.

\section{Conclusion}
Deep learning provides a rapid, computationally efficient triage surrogate for satellite conjunction risk analysis, delivering sub-millisecond classification and substantial speedups over classical Monte Carlo simulation.

\begin{thebibliography}{00}
\bibitem{b1} A. Vaswani et al., ``Attention is all you need,'' in \textit{NeurIPS}, 2017.
\bibitem{b2} Y. Sun et al., ``Retentive Network: A successor to Transformer for large language models,'' \textit{arXiv:2307.08621}, 2023.
\bibitem{b3} H. Touvron et al., ``LLaMA: Open and efficient foundation language models,'' \textit{arXiv:2302.13971}, 2023.
\bibitem{b4} T.-Y. Lin et al., ``Focal loss for dense object detection,'' in \textit{IEEE ICCV}, 2017.
\bibitem{b5} J. L. Foster and H. S. Estes, ``A Parametric Analysis of Orbital Debris Collision Probability,'' NASA JSC-25898, 1992.
\bibitem{b6} M. R. Akella and K. T. Alfriend, ``Probability of collision between space objects,'' \textit{JGCD}, vol. 23, no. 4, 2000.
\bibitem{b7} K. Chan, ``Collision probability for satellites in close encounters,'' \textit{JGCD}, vol. 20, no. 5, 1997.
\bibitem{b8} R. P. Patera, ``General method for calculating satellite collision probability,'' \textit{JGCD}, vol. 24, no. 4, 2001.
\bibitem{b9} K. T. Alfriend et al., ``Probability of collision error analysis,'' \textit{Space Debris}, vol. 1, no. 1, 1999.
\bibitem{b10} V. T. Coppola, ``Including velocity uncertainty in the probability of additional conjunctions,'' \textit{JGCD}, vol. 35, no. 5, 2012.
\bibitem{b11} D. T. Hall et al., ``Time-evolution of conjunction assessment risk metrics,'' in \textit{AAS/AIAA}, AAS 17-268, 2017.
\bibitem{b12} A. S. Mashiku and M. D. Hejduk, ``Recommended methods for setting conjunction assessment risk thresholds,'' \textit{JSSE}, vol. 6, no. 4, 2019.
\bibitem{b13} M. D. Hejduk and D. E. Snow, ``The state of conjunction assessment risk analysis,'' in \textit{AMOS}, 2019.
\bibitem{b14} T. Uriot et al., ``Spacecraft collision avoidance challenge,'' in \textit{NeurIPS Competition Track}, 2020.
\bibitem{b15} G. Acciarini et al., ``Spacecraft collision risk assessment with Gaussian processes,'' \textit{IEEE TAES}, vol. 57, no. 1, 2020.
\bibitem{b16} V. M. Pinto et al., ``Machine learning for spacecraft collision avoidance: A review,'' \textit{Acta Astronautica}, vol. 175, 2020.
\bibitem{b17} E. Serra et al., ``Deep learning for spacecraft collision avoidance maneuver planning,'' in \textit{AAS/AIAA}, 2021.
\bibitem{b18} D. Yepes et al., ``Automated satellite conjunction assessment using machine learning,'' \textit{JSSE}, vol. 9, no. 4, 2022.
\bibitem{b19} H. Peng and X. Bai, ``Artificial neural network for satellite collision probability calculation,'' \textit{Acta Astronautica}, vol. 151, 2018.
\bibitem{b20} T. G. Dietterich, ``Approximate statistical tests for comparing supervised classification learning algorithms,'' \textit{Neural Computation}, vol. 10, no. 7, 1998.
\end{thebibliography}

\end{document}
